% Complete random-series proofs, plus the source's explicitly conditional LF route.
\section{Independent random series: proofs}
\label{sec:annex-p15-independent-random-series}

The definitions and theorem statements are in
Chapter~\ref{sec:p15-independent-random-series}. The maximal and
variance-summability arguments follow the scope of Durrett's
Theorems 2.5.5--2.5.6. The elementary bounded converse supplies a
complete necessity proof for Theorem 2.5.8 without a central-limit
premise. The distinct source argument from Example 3.4.12 is preserved
in its own conditional item
\cite[pp.~84--85, 148--151]{DurrettPTEAuthor2019}.

\begin{annexitem}{annex:p15.series-event-calculus}{thm:p15-series-event-calculus}{Measurability, the tail property and finite changes}
\begin{proof}
Every partial sum is real measurable. Each set
\(\{|S_m-S_n|\leq1/r\}\) is consequently measurable, and the
displayed countable formula is a measurable event. If a sequence
converges in \(\mathbb R\), it is Cauchy and belongs to this event.
Conversely, for each \(\varepsilon>0\) choose an integer \(r\)
with \(1/r<\varepsilon\). The formula then gives an index after
which all pairwise differences are less than \(\varepsilon\).
Completeness of \(\mathbb R\) gives a finite limit. Countable
suprema of the measurable differences show that \(w_M\) is
measurable, possibly infinite. The sequence \(w_M\) decreases,
and the same Cauchy condition is exactly \(w_M\to0\).

Fix \(k\geq1\). Removing the first \(k-1\) summands subtracts
the finite number \(S_{k-1}\) from every subsequent partial sum.
It therefore preserves finite convergence. Apply the countable Cauchy
description to the partial sums starting at \(X_k\): all of its
events belong to \(\sigma(X_k,X_{k+1},\ldots)\). Since \(k\)
is arbitrary, \(C_X\) is a tail event. For independent summands,
Theorem~\ref{thm:p15-kolmogorov-tail-zero-one} gives probability
zero or one. This conclusion concerns the event of convergence; the
value of the sum need not be a tail random variable.

At an outcome where \(S_n\) has a finite limit,
\(X_n=S_n-S_{n-1}\to0\). If two summand sequences differ at only
finitely many indices at that outcome, their partial sums differ by
a fixed finite constant after the largest changed index. Their series
therefore converge together. The two absolute-value series differ by
only finitely many finite terms, which also preserves absolute
convergence. The index of the last change may depend on the outcome;
these assertions are pointwise.
\end{proof}
\end{annexitem}

\begin{annexitem}{annex:p15.kolmogorov-maximal}{thm:p15-kolmogorov-maximal}{The first-crossing decomposition}
\begin{proof}
For \(1\leq k\leq n\), define the first-crossing event
\[
 A_k=\{|S_k|\geq x\}\cap
                  \bigcap_{j=1}^{k-1}\{|S_j|<x\}.
\]
These events are disjoint, their union is
\(\{\max_{1\leq k\leq n}|S_k|\geq x\}\), and
\(A_k\in\sigma(X_1,\ldots,X_k)\).
The future increment \(T_k=S_n-S_k\) is independent of that
sigma-algebra by Theorem~\ref{thm:p15-independent-blocks}, and
\(\mathbb ET_k=0\). Both \(S_k\) and \(T_k\) have finite
second moments. The random variable \(S_k\mathbf1_{A_k}\) is
integrable and is independent of \(T_k\); their product is
integrable and has expectation zero by
Theorem~\ref{thm:p15-independent-expectations}. For \(k=n\),
the same conclusion holds with \(T_n=0\). Thus
\[
 \begin{aligned}
 \mathbb ES_n^2
 &\geq\sum_{k=1}^n\mathbb E[S_n^2\mathbf1_{A_k}]\\
 &=\sum_{k=1}^n\mathbb E[S_k^2\mathbf1_{A_k}]
   +2\sum_{k=1}^n\mathbb E[S_kT_k\mathbf1_{A_k}]
   +\sum_{k=1}^n\mathbb E[T_k^2\mathbf1_{A_k}]\\
 &\geq x^2\sum_{k=1}^n\mathbb P(A_k).
 \end{aligned}
\]
The first inequality discards the nonnegative contribution outside
the disjoint union, and the last discards nonnegative squared future
increments. Centering gives \(\mathbb ES_n^2=\operatorname{Var}(S_n)\).
Variance addition for independent summands
(Theorem~\ref{thm:p15-independent-moments}) identifies this with
\(\sum_{j=1}^n\operatorname{Var}(X_j)\). Division by \(x^2>0\)
proves the claim.
\end{proof}
\end{annexitem}

\begin{annexitem}{annex:p15.variance-series-convergence}{thm:p15-variance-series-convergence}{From maximal control to almost-sure and mean-square convergence}
\begin{proof}
For integers \(N>M\geq0\), apply
Theorem~\ref{thm:p15-kolmogorov-maximal} to
\(X_{M+1},\ldots,X_N\). For \(\varepsilon>0\),
\[
 \mathbb P\!\left(\max_{M\leq m\leq N}|S_m-S_M|>\varepsilon\right)
       \leq\varepsilon^{-2}\sum_{j=M+1}^N\operatorname{Var}(X_j).
\]
The strict-threshold event is contained in the corresponding
non-strict event in that theorem. Letting \(N\to\infty\) and
using continuity from below gives
\[
 \mathbb P\!\left(\sup_{m\geq M}|S_m-S_M|>\varepsilon\right)
       \leq\varepsilon^{-2}\sum_{j>M}\operatorname{Var}(X_j).
\]
For every pair \(m,n\geq M\), the triangle inequality bounds
\(|S_m-S_n|\) by \(2\sup_{j\geq M}|S_j-S_M|\). Hence
\[
 \mathbb P(w_M>2\varepsilon)
       \leq\varepsilon^{-2}\sum_{j>M}\operatorname{Var}(X_j)
          \longrightarrow0.
\]
For any fixed \(r\geq1\), the events \(\{w_M>2/r\}\) decrease,
so their intersection has probability zero. On the complement of the
union of these null intersections, each threshold \(2/r\) is
eventually an upper bound for \(w_M\). Thus \(w_M\to0\) on a
single probability-one event. The measurable Cauchy criterion gives
finite convergence there.

Write \(S_\infty=\lim_nS_n\) on \(C_X\), and set it to zero
on \(C_X^c\). This is real measurable: the extended measurable
\(\limsup_nS_n\) equals a finite number on the measurable set
\(C_X\), and measurable restriction followed by the specified
constant extension gives the claimed variable. For fixed \(M\),
Fatou's lemma applied to \((S_N-S_M)^2\) as \(N\to\infty\)
gives
\[
 \mathbb E(S_\infty-S_M)^2
       \leq\liminf_{N\to\infty}\mathbb E(S_N-S_M)^2
       =\sum_{j>M}\operatorname{Var}(X_j).
\]
Here the almost-sure limit suffices for the integral inequality;
changing values on the null exceptional set does not change an
integral. Taking \(M=0\) proves \(S_\infty\in L^2\), and
letting \(M\to\infty\) proves mean-square convergence. By
Cauchy--Schwarz on the probability space,
\[
 |\mathbb ES_\infty-\mathbb ES_M|
       \leq\mathbb E|S_\infty-S_M|
       \leq\sqrt{\mathbb E(S_\infty-S_M)^2}\longrightarrow0,
\]
so \(\mathbb ES_\infty=0\).

To obtain equality in the tail second moment, fix \(M\) and put
\(B_N=S_N-S_M\), \(B=S_\infty-S_M\). We already have
\(\mathbb E(B_N-B)^2\to0\), and both \(\mathbb EB_N^2\)
and \(\mathbb EB^2\) are bounded by the finite variance tail.
The pointwise bound \((B_N+B)^2\leq2B_N^2+2B^2\) and
Cauchy--Schwarz give
\[
 |\mathbb EB_N^2-\mathbb EB^2|
 \leq\sqrt{\mathbb E(B_N-B)^2}\,
          \sqrt{\mathbb E(B_N+B)^2}\longrightarrow0.
\]
Finite variance addition now yields
\(\mathbb EB^2=\sum_{j>M}\operatorname{Var}(X_j)\).
The case \(M=0\), together with the zero mean, proves the stated
variance of the limit. No completeness theorem for an \(L^2\)
space is needed for this construction.
\end{proof}
\end{annexitem}

\begin{annexitem}{annex:p15.bounded-fourth-moment}{thm:p15-bounded-fourth-moment}{A fourth-moment lower bound and the centered bounded converse}
\begin{proof}
All products below are integrable because the summands are bounded.
In the expansion of \(W^4\), any product in which some index occurs
exactly once has expectation zero: independence factors out the mean
of that centered factor. The surviving products have either one index
four times, or two different indices twice each. For a fixed pair of
different indices there are \(\binom42=6\) placements. Therefore
\[
 \mathbb EW^4
   =\sum_j\mathbb EU_j^4
       +6\sum_{i<j}\mathbb EU_i^2\,\mathbb EU_j^2
   \leq M^2v+3v^2.
\]
The inequality uses \(U_j^4\leq M^2U_j^2\) almost surely and
\(2\sum_{i<j}a_ia_j\leq(\sum_ja_j)^2\) for
\(a_j=\mathbb EU_j^2\geq0\). Similarly, expanding the square
and eliminating cross terms gives \(\mathbb EW^2=v\).

Suppose \(v>0\) and set \(D=\{W^2>v/2\}\). Splitting the
second moment at this event and using Cauchy--Schwarz gives
\[
 v=\mathbb E[W^2\mathbf1_{D^c}]+
                     \mathbb E[W^2\mathbf1_D]
   \leq v/2+\sqrt{\mathbb EW^4\,\mathbb P(D)}.
\]
The fourth moment is positive, since a zero fourth moment would force
\(W=0\) almost surely and contradict \(v>0\). Squaring and
rearranging the nonnegative inequality gives
\[
 \mathbb P(D)\geq\frac{v^2}{4\mathbb EW^4}
                  \geq\frac{v}{4(M^2+3v)}.
\]
If \(v\geq M^2+1\), then \(M^2\leq v\) and \(v\geq1\),
so this probability is at least \(1/16\), while
\(\sqrt{v/2}\geq1/\sqrt2>1/2\). Thus whenever
\(v\geq M^2+1\), we have \(\mathbb P(|W|>1/2)\geq1/16\).

Now let \((U_j)_{j\geq1}\) be the uniformly bounded independent
centered sequence in the converse, and denote its partial sums by
\(T_n\). Almost-sure convergence implies that its oscillations
\(q_m=\sup_{i,j\geq m}|T_i-T_j|\) tend to zero almost surely.
The events \(\{q_m>1/2\}\) decrease to a null event, so
\(\mathbb P(q_m>1/2)\to0\).
If \(\sum_j\operatorname{Var}(U_j)=\infty\), every variance
tail is infinite. For each \(m\), choose a finite \(N>m\) with
\(\sum_{j=m+1}^N\operatorname{Var}(U_j)\geq M^2+1\).
Apply the just-proved finite-sum bound to this block:
\[
 1/16\leq\mathbb P(|T_N-T_m|>1/2)
                  \leq\mathbb P(q_m>1/2).
\]
This contradicts the limit as \(m\to\infty\). The variances
must therefore be summable.
\end{proof}
\end{annexitem}

\begin{annexitem}{annex:p15.series-independent-copy}{thm:p15-series-independent-copy}{Constructing and using the copied sequence}
\begin{proof}
The product probability measure exists by
Theorem~\ref{thm:p15-product-existence}; each factor is finite and
therefore sigma-finite. Coordinate maps on the product are measurable,
so all copied variables and their differences are measurable.
For finite index sets \(I,J\) and Borel sets \(B_i,D_j\), the
event imposing \(\widehat Y_i\in B_i\) for \(i\in I\) and
\(\widetilde Y_j\in D_j\) for \(j\in J\) is the rectangle
\[
 \left(\bigcap_{i\in I}\{Y_i\in B_i\}\right)
     \times\left(\bigcap_{j\in J}\{Y_j\in D_j\}\right).
\]
Its product probability is
\(\prod_{i\in I}\mathbb P(Y_i\in B_i)
  \prod_{j\in J}\mathbb P(Y_j\in D_j)\).
This is exactly finite-subfamily independence of the combined copied
family, even when an index occurs once in each copy.
The pairs \((\widehat Y_n,\widetilde Y_n)\) use disjoint members
of that family. Theorem~\ref{thm:p15-independent-blocks} and
measurability of subtraction imply independence of the \(Z_n\).

Each copy has the original law, hence
\(\mathbb EZ_n=\mathbb EY_n-\mathbb EY_n=0\).
The almost-sure bound is \(2A\). These bounds can be used
simultaneously for all indices because a countable union of the
exceptional null sets is null. Independence of the two copies gives
\(\operatorname{Var}(Z_n)=2\operatorname{Var}(Y_n)\).
The coordinate-swap map \((\omega,\omega')\mapsto(\omega',\omega)\)
is measurable and preserves \(\mathbb P\otimes\mathbb P\): its
pushforward agrees with that probability measure on every measurable
rectangle, and uniqueness on the generating rectangles gives equality
on the product sigma-algebra. Swapping changes \(Z_n\) to \(-Z_n\),
which proves symmetry of its law.

If the original convergence event \(C_Y\) has probability one,
then \(C_Y\times C_Y\) has product probability one. On this
rectangle, the partial sums of \(Z_n\) are differences of two
convergent real partial-sum sequences. They converge to the difference
of the two finite limits. This proves the assertion without requiring
an independent copy to exist on the original space.
\end{proof}
\end{annexitem}

\begin{annexitem}{annex:p15.bounded-series-criterion}{thm:p15-bounded-series-criterion}{The two bounded-series directions}
\begin{proof}
Suppose the two numerical conditions hold. The variables
\(Y_n-\mathbb EY_n\) are independent, centered and have summable
variances, so Theorem~\ref{thm:p15-variance-series-convergence}
gives almost-sure convergence of their series. Adding the convergent
numerical series \(\sum_n\mathbb EY_n\) proves almost-sure
convergence of \(\sum_nY_n\).

Conversely, suppose \(\sum_nY_n\) converges almost surely.
Construct the differences \(Z_n\) on the product space by
Lemma~\ref{thm:p15-series-independent-copy}. Their series converges
almost surely, and they are centered, independent and uniformly
bounded by \(2A\). The centered bounded converse in
Lemma~\ref{thm:p15-bounded-fourth-moment} yields
\[
 2\sum_n\operatorname{Var}(Y_n)
                  =\sum_n\operatorname{Var}(Z_n)<\infty.
\]
Return to the original probability space. The variance-summability
theorem now proves almost-sure convergence of
\(\sum_n(Y_n-\mathbb EY_n)\). Intersect this probability-one
event with the probability-one convergence event for \(\sum_nY_n\).
On the intersection,
\[
 \sum_{j=1}^n\mathbb EY_j
   =\sum_{j=1}^nY_j-\sum_{j=1}^n(Y_j-\mathbb EY_j)
\]
converges to a finite number. The expression on the left is a
deterministic numerical sequence; the nonempty probability-one
intersection supplies an outcome witnessing its convergence.
Thus the numerical mean series converges. This argument uses neither
the three-series theorem nor a central-limit theorem.
\end{proof}
\end{annexitem}

\begin{annexitem}{annex:p15.three-series}{thm:p15-three-series}{Large jumps, drift and centered fluctuations}
\begin{proof}
Fix \(A>0\) and abbreviate \(Y_n=Y_n^{(A)}\). The map
\(x\mapsto x\mathbf1_{\{|x|\leq A\}}\) is Borel measurable.
Therefore the \(Y_n\) are independent measurable functions of the
\(X_n\), and satisfy \(|Y_n|\leq A\) everywhere. In particular,
each mean and variance in the criterion is finite.

For sufficiency, the second and third conditions and
Theorem~\ref{thm:p15-bounded-series-criterion} give almost-sure
convergence of \(\sum_nY_n\). Equivalently, apply
Theorem~\ref{thm:p15-variance-series-convergence} to
\(Y_n-\mathbb EY_n\) and add the convergent mean series.
The first Borel--Cantelli lemma applied to \(\{|X_n|>A\}\)
shows that \(X_n\neq Y_n\) for only finitely many indices at
almost every outcome. The finite-change property in
Theorem~\ref{thm:p15-series-event-calculus} then gives almost-sure
convergence of \(\sum_nX_n\).

For necessity, almost-sure finite convergence forces \(X_n\to0\)
almost surely. The events \(\{|X_n|>A\}\) are independent,
since they are separate coordinate events of the independent sequence.
If their probabilities had infinite sum, the second Borel--Cantelli
lemma (Theorem~\ref{thm:p15-borel-cantelli-second}) would make
\(|X_n|>A\) occur infinitely often almost surely. This contradicts
\(X_n\to0\), since \(A>0\). Thus the first condition holds.
The first Borel--Cantelli lemma now gives only finitely many changed
summands almost surely, so \(\sum_nY_n\) converges almost surely.
The necessity direction of
Theorem~\ref{thm:p15-bounded-series-criterion} supplies the second
and third conditions. This proves both directions at the fixed cutoff.

The necessity proof works for every \(A>0\), while sufficiency at
one such cutoff already proves convergence. This gives the equivalence
between one-cutoff and every-cutoff satisfaction. If any condition
fails at a fixed cutoff, almost-sure convergence is impossible.
The event of convergence is a tail event of an independent sequence,
so its probability is zero or one. Excluding one therefore proves
that its probability is zero.
\end{proof}
\end{annexitem}

\begin{annexitem}{annex:p15.three-series-source-route}{thm:p15-three-series-source-route}{The complete source necessity route using Lindeberg--Feller}
\begin{proof}
We use precisely \(\mathrm{LF}\) from
Definition~\ref{def:p15-series-lindeberg-premise}, now proved in
\cref{thm:p15-clt-lindeberg-feller}. The following
argument preserves Durrett's Example 3.4.12; the distribution-function
step invoked there by an exercise is supplied explicitly here.

Let \(\sum_nX_n\) converge almost surely and fix \(A>0\).
As in the direct necessity proof, the second Borel--Cantelli lemma
and \(X_n\to0\) prove
\(\sum_n\mathbb P(|X_n|>A)<\infty\). The first lemma then
implies that \(Y_n=X_n\mathbf1_{\{|X_n|\leq A\}}\) differs
from \(X_n\) finitely often almost surely. Consequently its raw
partial sums \(R_n=\sum_{m=1}^nY_m\) converge to a finite
random limit almost surely.

Suppose for a contradiction that
\(c_n=\sum_{m=1}^n\operatorname{Var}(Y_m)\to\infty\).
For each row with \(c_n>0\), put
\[
 U_{n,m}=\frac{Y_m-\mathbb EY_m}{\sqrt{c_n}}
       \quad(1\leq m\leq n),\qquad
 Z_n=\sum_{m=1}^nU_{n,m},\qquad
 V_n=\frac{R_n}{\sqrt{c_n}}.
\]
For any initial row with \(c_n=0\), define all its \(U_{n,m}\),
and \(Z_n,V_n\), to be zero. Since \(c_n\to\infty\), only
finitely many rows have zero \(c_n\), so these choices have no
effect on the limits below.

Each nonzero row is independent and centered, and
\[
 \sum_{m=1}^n\mathbb EU_{n,m}^2=1,\qquad
 |U_{n,m}|\leq\frac{2A}{\sqrt{c_n}}.
\]
For a given \(\varepsilon>0\), the last bound is at most
\(\varepsilon\) for all sufficiently large \(n\) and every
index \(m\) in that row. The whole Lindeberg sum is then zero.
Thus the two premises of \(\mathrm{LF}\) hold with \(\sigma^2=1\).
It follows that
\[
 \mathbb P(Z_n\leq t)\longrightarrow\Phi(t)
                              \quad(t\in\mathbb R).
\]
Because \(R_n\) converges to a finite value almost surely and
\(\sqrt{c_n}\to\infty\), we have \(V_n\to0\) almost surely.
For every \(\delta>0\), bounded convergence applied to
\(\mathbf1_{\{|V_n|>\delta\}}\) gives
\(\mathbb P(|V_n|>\delta)\to0\).

Here is the required perturbation argument. For any real \(t\)
and \(\delta>0\), elementary event inclusions give
\[
 \begin{aligned}
 \mathbb P(Z_n\leq t-\delta)-\mathbb P(|V_n|>\delta)
  &\leq\mathbb P(Z_n-V_n\leq t)\\
  &\leq\mathbb P(Z_n\leq t+\delta)
                           +\mathbb P(|V_n|>\delta).
 \end{aligned}
\]
Indeed, \(Z_n\leq t-\delta\) and \(|V_n|\leq\delta\)
imply \(Z_n-V_n\leq t\); conversely,
\(Z_n-V_n\leq t\) and \(|V_n|\leq\delta\) imply
\(Z_n\leq t+\delta\). Taking lower and upper limits bounds
the distribution-function limit between \(\Phi(t-\delta)\)
and \(\Phi(t+\delta)\). The standard normal law has a density,
so it gives zero probability to singletons and its distribution
function is continuous. Letting \(\delta\downarrow0\) proves
\[
 \mathbb P(Z_n-V_n\leq t)\longrightarrow\Phi(t).
\]

But for every sufficiently large \(n\),
\[
 Z_n-V_n=-\frac{\sum_{m=1}^n\mathbb EY_m}{\sqrt{c_n}}
                   =d_n
\]
is a deterministic real number. At \(t=0\), its distribution
function is \(\mathbf1_{\{d_n\leq0\}}\), which is either
zero or one. It cannot converge to \(\Phi(0)=1/2\): every
term has distance \(1/2\) from that number. The identity
\(\Phi(0)=1/2\) follows from the even standard normal density,
its normalization and its zero mass at zero. This is the desired
contradiction. Thus \(\sum_n\operatorname{Var}(Y_n)<\infty\).

Theorem~\ref{thm:p15-variance-series-convergence} now proves
almost-sure convergence of \(\sum_n(Y_n-\mathbb EY_n)\).
Subtracting its partial sums from the almost-surely convergent
\(R_n\), on the intersection of the two probability-one events,
proves convergence of the deterministic series \(\sum_n\mathbb EY_n\).
This completes the source's alternate necessity route.

The theorem premise \(\mathrm{LF}\), namely Durrett's Theorem 3.4.10,
is supplied by the characteristic-function construction in
\cref{thm:p15-clt-lindeberg-feller}, with the separate moment proof
in \cref{thm:p15-clt-moment-array-proof}.
The elementary necessity proof in
Annex~\ref{annex:p15.three-series} already establishes the
unconditional three-series theorem independently of this premise.
\end{proof}
\end{annexitem}

\begin{annexitem}{annex:p15.signed-power-series}{thm:p15-signed-power-series}{Variance summability and the two power thresholds}
\begin{proof}
First record the numerical power-series criterion. For \(p>0\)
and each integer \(k\geq0\), there are \(2^k\) terms in the
block \(2^k\leq n<2^{k+1}\), and
\[
 2^{-p}2^{k(1-p)}
       \leq\sum_{n=2^k}^{2^{k+1}-1}n^{-p}
       \leq2^{k(1-p)}.
\]
If \(p>1\), the upper bounds form a convergent geometric series.
If \(0<p\leq1\), the lower bounds do not tend to zero, so the
sum of blocks diverges. For \(p\leq0\), the terms \(n^{-p}\)
themselves do not tend to zero. Thus \(\sum_nn^{-p}\) is finite
exactly for \(p>1\).

For \(X_n=\varepsilon_nn^{-\alpha}\), the variables are
independent measurable transforms of the signs, and
\(\mathbb EX_n=0\), \(\operatorname{Var}(X_n)=n^{-2\alpha}\).
If \(\alpha>1/2\), the variance sum is finite by the just-proved
criterion, and Theorem~\ref{thm:p15-variance-series-convergence}
gives almost-sure convergence.
If \(0<\alpha\leq1/2\), all \(|X_n|\leq1\), so at cutoff
one the truncated variables are exactly \(X_n\); the variance
series diverges. Theorem~\ref{thm:p15-three-series} gives
probability zero for finite convergence. If \(\alpha\leq0\),
\(|X_n|=n^{-\alpha}\) does not tend to zero at any outcome, so
finite convergence is impossible there as well.

At every outcome, \(\sum_n|X_n|=\sum_nn^{-\alpha}\).
The numerical criterion makes this finite exactly when \(\alpha>1\).
For \(1/2<\alpha\leq1\), finite convergence holds almost surely
but absolute convergence fails at every outcome, giving the asserted
conditional convergence.
\end{proof}
\end{annexitem}

\begin{annexitem}{annex:p15.series-boundary-examples}{thm:p15-series-boundary-examples}{Computing each hypothesis boundary}
\begin{proof}
Deterministic random variables are independent: any event constraining
a constant variable has probability zero or one, and finite
intersection probabilities factor. For \(X_n=(-1)^{n+1}/n\),
all values have absolute value at most one, all variances are zero,
and the means are the same alternating numbers. To check convergence
without an additional series theorem, let \(a_n\) be their partial
sums. The even subsequence increases because
\(a_{2k+2}-a_{2k}=1/(2k+1)-1/(2k+2)>0\).
The odd subsequence decreases because
\(a_{2k+3}-a_{2k+1}=-1/(2k+2)+1/(2k+3)<0\).
Moreover \(0\leq a_{2k}<a_{2k+1}\leq1\), and their gap is
\(1/(2k+1)\to0\). The two bounded monotone subsequences
therefore have the same finite limit. The harmonic absolute-value
series diverges by the dyadic-block calculation in
Annex~\ref{annex:p15.signed-power-series}. This verifies the
conditional mean-series example.

For the first example separating the three conditions, the stated
three masses are nonnegative and sum to one, including \(n=1\),
where the zero mass is zero. Every nonzero value has absolute value
\(n+1>1\), so \(Y_n^{(1)}=0\). The mean and variance series
are therefore zero. The large-jump probability is \(1/n\), whose
sum diverges. Independence and the second Borel--Cantelli lemma make
these nonzero values occur infinitely often almost surely. The terms
then fail to tend to zero, proving almost-sure divergence directly.
For deterministic \(X_n=1/n\), the retained mean series is the
divergent harmonic series, while no large jumps occur and every
variance is zero. The partial sums diverge to positive infinity.
For independent \(X_n=\varepsilon_n/\sqrt n\), there are no
large jumps at cutoff one, every mean is zero, and the retained
variance is \(1/n\). The signed-power theorem at \(\alpha=1/2\)
proves almost-sure divergence. Thus each separate condition is needed.

For the rare large summands in the third example,
\[
 \mathbb P(X_n\ne0)=2^{-n},\qquad
 \mathbb EX_n=0,\qquad
 \mathbb EX_n^2=(2^n)^2\,2^{-n}=2^n.
\]
Each mean and variance exists because each individual law has finite
support. The sum of nonzero probabilities is finite. The first
Borel--Cantelli lemma implies that only finitely many summands are
nonzero almost surely. Each is a finite real number, so their series
converges absolutely on that event. Its original variance series
diverges, as the displayed calculation shows. At cutoff one all
nonzero values are deleted; the retained means and variances are
zero and the large-jump probabilities have finite sum, consistently
with the three-series criterion.

Finally, for the shared-sign sequence \(X_n=\varepsilon/n\),
\(\mathbb EX_n=0\) and \(\operatorname{Var}(X_n)=1/n^2\),
so the variance sum is finite. However, its partial sums are
\(\varepsilon\sum_{j=1}^n1/j\), which diverge to positive
infinity or negative infinity according to the sign. The variables
are dependent; for example the events that \(X_1\) and \(X_2\)
are positive coincide and have probability \(1/2\), while their
probability product is \(1/4\).
For the centered dependent pair \(\varepsilon,-\varepsilon\),
\(S_1=\varepsilon\) and \(S_2=0\). At \(x=1\), the left
side of the claimed maximal inequality would be one, while the
terminal variance on the right is zero. For the uncentered variable
\(X_1=1\), the same threshold gives probability one and variance
zero. These computations prove the two maximal-inequality failures.
The independent examples exist on countable product probability
spaces by Theorem~\ref{thm:p15-independent-sequence-existence};
the dependent examples use just one two-point probability space.
\end{proof}
\end{annexitem}
